\documentclass[../paper.tex]{subfiles}

\begin{document}

\subfile{../figures/comparison.tex}

The overview and taxonomy of the clustering algorithms included in this survey follow in Section~\ref{sec:methods}. First, we provide a brief overview of relevant objective functions, algorithms, and survey papers. The main motivation for this survey, and reason for focusing primarily on modularity optimization, is highlighted in Figure~\ref{fig:comparisons}. Despite high interest from multiple scientific communities and the rapid development of new methods, comparisons between methods remain fragmented. This is a problem for both practitioners who want the best algorithm and for future researchers who need to decide which baseline algorithms to use.

While modularity has attracted significant interest, it is not without its flaws. The most well-known problem is the resolution limit~\cite{resolution2007}. In extreme cases, this can cause small communities with few edges connecting them to appear as one larger community. As a result, several algorithms provide a \textsc{resolution} parameter to detect smaller communities in large graphs~\cite{louvain2008, leiden2019}.

There have also been entirely new objective functions with implemented algorithms available. One such example is the \textit{Weighted Community Cluster} (WCC), which maximizes the number of triangles within each community~\cite{wcc2012}. The \textsc{Scalable Community Detection} algorithm was later developed for the WCC objective~\cite{scd2014}. Another objective function is the \textit{Map Equation} that measures the expected description length of a random walk~\cite{infomap2008}. Several algorithms have been developed for the map equation objective, such as \textsc{Infomap}~\cite{infomap2008}, \textsc{Infoflow}~\cite{infoflow2019}, and the very recent learning-based \textsc{Neuromap}~\cite{neuromap2024}.

There have also been similar experimental surveys looking at modularity. A study by Aref et al.~\cite{survey2023} found that modularity-optimizing heuristics rarely find optimal solutions. That claim was later disputed by Ans{\'o}tegui et al.~\cite{maxsat2025} who found that the \textsc{VieClus} heuristic routinely computes optimal solutions on small to medium-sized graphs.

% Maybe mention editing distance

\ifSubfilesClassLoaded{%
    \bibliography{../paper.bib}%
}{}

\end{document}